\documentclass[12pt, a4paper]{article}
\usepackage{amsmath, amsfonts}
\usepackage[english]{babel}

\begin{document}

\section*{Problem Statement}

Let $f : \mathbb{Z} : \mathbb{Z}$ such that
\[
\forall x,y\in \mathbb Z:\quad f(xy-1) + f(x+y) = f(x)f(y)+2.
\]
Find $f(45)$.
\section*{Solution}

Whenever we see identities involding sums and products of function argument,
it is useful to start with values like $0$ and $\pm1$.
Here, if we denote $z = f(0)$, $m = f(-1)$, and $p=f(1)$, we can write (take $y=1$)
\[
f(x-1) + f(x+1) = p f(x) + 2;\]
This recurrent formula will the basis of further analysis.

On the other hand, let us take $y=0$; this results in
\[
m + f(x) = zf(x) + 2.
\]
In particular, for $x=0$ we get
\[
m + z = z^2 + 2
\]
and for $x=-1$ we get
\[2m=mz+2.\]
The latter equation, rewritten as $m(2-z) = 2$, andmits only four integer solutions:
$(m=2, z=1)$, $(m=-2, z=3)$, $(m=1, z=0)$, and $(m=-1, z=4)$.
Only the first one satisfies the equation $m+z=z^2+2$ (todo: label instead rewrite).

Now, if we manage to find $p$, any value of $f$ can be found.
In order to do so, we will take a positive integer number with two different decompositions into a sum of two positive numbers.
Obviously, $4$ is smallest such number: $4 = 1+3=2+2$, so we write
\[f(2\cdot2-1)+f(2+2) = f^2(2)+2,\quad f(1\cdot3-1)+f(1+3)=pf(3)+2.\]
In other words
\[
f(4) = f^2(2) + 2 - f(3) = pf(3)+2-f(2).
\]
On the other hand, thanks to our recurrence relation,
\[f(2) = p^2 +2 - z = p^2+1,\quad f(3) = pf(2)+ 2 - p = p^3+2.\]
Substituting these into the previous equation, we obtain
\[
p^4+2p^2+1 +2- p^3-2 = p^4 + 2p +2 - p^2 -1,
\]
or
\[
p^3 - 3p^2 + 2p   = 0.
\]
This equation has three solutions, $p=0,1,2$ (we, of course hoped for a single integer solution).
Let's write also \[f(4) = p^4 + 2p + 2 - p^2 -1 = p^4 - p^2 + 2p +1.\]


Let us study $5 = 2+3 = 1+4$:
\[f(2\cdot 3 - 1) + f (2+3) = f(2) f(3)+2,\quad f(1\cdot4-1)+f(1+4)=pf(4)+2.\]
or
\[
2f(5) = f(2)f(3)+2,\quad f(5) = p f(4)+2 - f(3).\]
Substituting the values of $f(2)$ and $f(3)$ would yield
\[
2f(5)= p^5 + p^3 + 2p^2 + 4, \quad 2f(5) = 2p^5-4p^3+4p^2+2p.
\]
Simplifying the difference of these two identities, we get
\[
p^5-5p^3+2p^2+2p-4 = 0.
\]
As we already know, the candidates for $p$ are $0,1,2$.
Clearly, only $p=2$ satisfies the equation.

After that, the recurrent relation becomes
\[f(x+1) = 2f(x) - f (x-1) + 2,\]
and its terms are (starting with $x=-1$) are
\[2, 1, 2, 5, 10, 17,\ldots\]

Let us now take $g(x) = f(x) - f(x-1)$, $g(0) = z - m = -1$ and rewrite
\[g(x+1) = g(x) + 2\]
with an evident solution
\[
g (x) = g(0) + 2x = 2 x - 1.
\]
Now we can write
\[
f(t) - f(t-1) = g(t) = 2 t - 1
\]
and sum over $t$:
\[
f(x) - f(0) = \sum_{t=1}^x (f(t) - f(t-1)) = \sum_{t=1}^x (2 t - 1) = x^2.
\]
Finally, $f(x) = f(0) + x^2 = 1 + x^2$.

Among other things, $f(45) = 2026$.

\end{document}
